\section{Step 29: Prepare a private record context}
Multiply the admitted frame count by eight and compare it with the native
directory byte capacity at context offset zero. Reject a short capacity before
reconstructing any scratch frames. Eight times at most 256 fits easily in a
64-bit register.

Fill \texttt{index\_capacity}. Initialize the private 48-byte record context
at \texttt{\$sp+128}: capacity 16\,384, expected sequence zero, the admitted
period, and three zero fields. Its sequence will advance as each nested record
call succeeds. This is distinct from the public index context in
\texttt{\$s3}. No successful record may publish index metadata through that
private pointer.

\paragraph{Observe} Prepare \texttt{s29-1} and stop before
\texttt{index\_capacity\_ready}. At \texttt{index\_capacity\_ready}, the required directory
extent in \texttt{\$t1} is 64. Private record-context fields at offsets 128,
136 and 144 are 16\,384, zero and four. The public directory is unchanged.

\Needspace{7\baselineskip}
\paragraph{Check} Which sequence belongs in the private context initially?
\begin{itemize}
  \item[A.] Zero.
  \item[B.] The number of frames.
  \item[C.] The file byte count.
\end{itemize}
